\section*{Chapter \ref{ch:b1:operational-amplifier} --- The Operational Amplifier}

\begin{solution}{exo:b1:operational-amplifier:1}
$G = 1 + 9 = 10$; $V_s = \qty{5.0}{V}$; saturation at $V_e = 15/10 =
\qty{1.5}{V}$.
\end{solution}

\begin{solution}{exo:b1:operational-amplifier:2}
$G = -10$; input \omterm{def:b1:sinusoidal-impedance:impedance}{impedance} $R_1 = \qty{10}{k\ohm}$; $i = 1.0/10^4 =
\qty{0.10}{mA}$; $V_s = \qty{-10}{V}$.
\end{solution}

\begin{solution}{exo:b1:operational-amplifier:3}
Without: divider $1/(1 + 10) = 0.091$, \qty{91}{mV}. With the follower:
\qty{1.0}{V} --- the follower draws nothing from the sensor and supplies
the \qty{1}{mA} the load needs.
\end{solution}

\begin{solution}{exo:b1:operational-amplifier:4}
$V_s = -R_f(V_1/R_1 + V_2/R_2) = -(V_1 + V_2/2) = -(1.0 + 1.0) = \qty{-2.0}{V}$.
\end{solution}

\begin{solution}{exo:b1:operational-amplifier:5}
$V_+ = V_2R_2/(R_1 + R_2)$; at $V_-$: $(V_1 - V_-)/R_1 = (V_- - V_s)/R_2$,
so $V_s = V_-(1 + R_2/R_1) - V_1R_2/R_1$; with $V_- = V_+$:
$V_s = V_2R_2/R_1 - V_1R_2/R_1$. Numbers: $10 \times 0.05 = \qty{0.50}{V}$;
the common \qty{1}{V} contributes nothing --- common-mode rejection.
\end{solution}

\begin{solution}{exo:b1:operational-amplifier:6}
$V_s = -(1/RC)\int V_e$: slope $-1000 \times 1.0 = \qty{-1000}{V/s}$;
saturation after $15/1000 = \qty{15}{ms}$. Square wave: triangle of
\omterm{def:b1:wave-propagation:signal}{amplitude} $ET/4RC = 1.0 \times 10^{-3}/(4 \times 10^{-3}) = \qty{0.25}{V}$
(plus whatever constant the initial charge left).
\end{solution}

\begin{solution}{exo:b1:operational-amplifier:7}
$\underline H = -(R'/R)/(1 + jR'C\omega)$: DC gain $-100$; cutoff
$1/(2\pi R'C) = \qty{1.6}{Hz}$; at \qty{1}{kHz}, $R'C\omega = 628$, $G =
100/628 = 0.16$ --- the integrator value $1/RC\omega$. It integrates for
$f \gg \qty{1.6}{Hz}$.
\end{solution}

\begin{solution}{exo:b1:operational-amplifier:8}
$R_2/R_1 = 10$: $R_2 = \qty{10}{k\ohm}$; $C = 1/(2\pi R_2f_c) = 1/(2\pi
\times 10^4 \times 500) = \qty{32}{nF}$.
\end{solution}

\begin{solution}{exo:b1:operational-amplifier:9}
With finite $\mu$: $V_s = \mu(V_e - \beta V_s)$, $\beta = 1/100$, so
$G = \mu/(1 + \mu\beta) = 10^5/(1 + 10^3) = 99.9$. Bandwidth $f_T/G =
\qty{10}{kHz}$; follower: \qty{1}{MHz}.
\end{solution}

\begin{solution}{exo:b1:operational-amplifier:10}
Each noise excursion across \qty{2.0}{V} flips the output: a burst of
rapid switching (``chatter'') around the crossing. Remedy: hysteresis
wider than the noise, i.e.\ a \omterm{prop:b1:operational-amplifier:schmitt}{Schmitt trigger}.
\end{solution}

\begin{solution}{exo:b1:operational-amplifier:11}
$V_+ = V_sR_1/(R_1 + R_2) = V_s/10$: thresholds $\pm\qty{1.5}{V}$. The
output is $+\qty{15}{V}$ while $V_e$ rises until \qty{1.5}{V}, flips to
$-\qty{15}{V}$, stays there until $V_e$ falls below \qty{-1.5}{V}, flips
back: a square wave of the triangle's frequency, its edges at the
threshold crossings.
\end{solution}

\begin{solution}{exo:b1:operational-amplifier:12}
With $V_s = +V_{\mathrm{sat}}$, $V_C$ charges from $-\beta V_{\mathrm{sat}}$
toward $+V_{\mathrm{sat}}$: $V_C = V_{\mathrm{sat}} - (1 + \beta)V_{\mathrm{sat}}
\eu^{-t/RC}$, reaching $+\beta V_{\mathrm{sat}}$ when $\eu^{-t/RC} =
(1 - \beta)/(1 + \beta)$: half-period $RC\ln\frac{1 + \beta}{1 - \beta}$,
hence $T$. $\beta = \tfrac12$: $T = 2RC\ln 3 = 2.2RC$; $RC = 1/(2.2 \times
10^3) = \qty{0.455}{ms}$, $R = \qty{4.6}{k\ohm}$. $V_s$: square wave
$\pm\qty{15}{V}$; $V_C$: exponential arcs between $\pm\qty{7.5}{V}$.
\end{solution}

\begin{solution}{pb:b1:operational-amplifier:1}
\textbf{1.} Two dividers: $u = E[\tfrac12 - (1 + \epsilon)/(2 + \epsilon)]
\approx -E\epsilon/4$ (\cref{prop:b1:dc-circuits:wheatstone}); wiring
fixes the sign.

\textbf{2.} $u = 5.00 \times 2.0\times10^{-3}/4 = \qty{2.5}{mV}$.

\textbf{3.} $V_A = E/2 = \qty{2.500}{V}$ always; $V_B = E(1 + \epsilon)/(2 +
\epsilon) \approx 2.500 + 0.0025 = \qty{2.5025}{V}$ at full scale. The
\qty{2.5}{V} shared by both is the common-mode \omterm{def:b1:dc-circuits:voltage}{voltage}.

\textbf{4.} Each midpoint: $1 \parallel 1 = \qty{500}{\ohm}$; between the
two midpoints, \qty{1.0}{k\ohm}.

\textbf{5.} $10/(10 + 1) = 0.91$: a $9\%$ loss.

\textbf{6.} Infinite input \omterm{def:b1:sinusoidal-impedance:impedance}{impedance} (no current from the bridge, so
the midpoints keep their open-circuit potentials) and zero output
\omterm{def:b1:sinusoidal-impedance:impedance}{impedance} (the next stage can draw what it needs).

\textbf{7.} $V_A$ and $V_B$ exactly; zero current.

\textbf{8.} As in \cref{prop:b1:operational-amplifier:sumdiff}: $V_+ =
V_BR_2/(R_1 + R_2)$, \omterm{thm:b1:dc-circuits:kirchhoff}{node law} at $V_-$, $V_- = V_+$.

\textbf{9.} $R_2 = \qty{1.0}{M\ohm}$.

\textbf{10.} $1 + R_2/R_1 = 10$: e.g.\ \qty{1}{k\ohm} and \qty{9}{k\ohm}.
Total gain $1000$; full scale $2.5 \times 10^{-3} \times 1000 = \qty{2.5}{V}$.

\textbf{11.} Matched: the common \qty{2.5}{V} cancels exactly (only
$V_B - V_A$ survives). A $1\%$ mismatch passes about $1\%$ of it: an
input-referred error of order \qty{25}{mV}, ten times the full-scale
\omterm{def:b1:wave-propagation:signal}{signal} --- the four \omterm{def:b1:dc-circuits:resistor}{resistors} must be matched to better than $0.01\%$
(or the common-mode \omterm{def:b1:dc-circuits:voltage}{voltage} removed first).

\textbf{12.} $1\ \unit{mV} \times 1000 = \qty{1}{V}$, forty percent of full
scale: unusable without offset trimming, a low-offset (chopper)
amplifier, or a zero calibration with no load on the girder.

\textbf{13.} Stage 1: $f_T/100 = \qty{10}{kHz}$; stage 2: \qty{100}{kHz};
far above \qty{5}{Hz}.

\textbf{14.} Required $2.5/10^{-2} = \qty{250}{V/s} = \qty{0.25}{mV/\micro s}$,
two thousand times below the \omterm{rem:b1:operational-amplifier:real}{slew rate}: no.

\textbf{15.} $C = 1/(2\pi R_2f_c) = 1/(2\pi \times 1.6\times10^4 \times 10)
= \qty{1.0}{\micro F}$.

\textbf{16.} $x = 5$: $G = 1/\sqrt{26}$, \qty{-14}{dB}.

\textbf{17.} $\varphi = -\arctan 0.5 = -27^\circ$; delay $\varphi/\omega =
\qty{15}{ms}$ --- nothing against a truck's passage.

\textbf{18.} No loading of the following stage (zero output \omterm{def:b1:sinusoidal-impedance:impedance}{impedance})
and gain available; also no attenuation of the wanted \omterm{def:b1:wave-propagation:signal}{signal} below
cutoff.

\textbf{19.} \qty{-28}{dB} (gains multiply, \omterm{def:b1:filters-transfer-functions:bode}{decibels} add), exact because
the first stage's output \omterm{def:b1:sinusoidal-impedance:impedance}{impedance} is zero: the second stage does not
load it.

\textbf{20.} $\epsilon = 1.6\times10^{-3} \to u = \qty{2.0}{mV} \to
\qty{2.0}{V}$; output $+V_{\mathrm{sat}}$ when the \omterm{def:b1:wave-propagation:signal}{signal} exceeds
\qty{2.0}{V}, $-V_{\mathrm{sat}}$ below (or the reverse, by wiring).

\textbf{21.} It chatters: the lamp flickers as noise crosses the
reference.

\textbf{22.} With $V_{\mathrm{ref}}$ through $R_2$ and the output through
$R_1$... in the standard form, the non-inverting input sits at
$V_{\mathrm{ref}}(1 - \beta) + \beta V_s$; taking the reference shifted
accordingly, the switching levels are $V_{\mathrm{ref}} \pm \beta V_{\mathrm{sat}}$.
$\beta V_{\mathrm{sat}} = \qty{0.10}{V}$: $\beta = 1/150$, e.g.\ $R_1 =
\qty{1.0}{k\ohm}$, $R_2 = \qty{149}{k\ohm}$.

\textbf{23.} \qty{2.1}{V} (trip) and \qty{1.9}{V} (release): strains
\num{1.68e-3} and \num{1.52e-3}.

\textbf{24.} $R = (15 - 2.0)/0.013 = \qty{1.0}{k\ohm}$. At
$-V_{\mathrm{sat}}$ the LED is reverse-biased by \qty{15}{V}, beyond its
rating: an ordinary diode in antiparallel (or in series) protects it.

\textbf{25.} Bridge (\qty{2.5}{mV} full scale) $\to$ two followers (no
loading) $\to$ \omterm{prop:b1:operational-amplifier:sumdiff}{difference amplifier} $\times100$ $\to$ non-inverting
$\times10$ $\to$ active low-pass \qty{10}{Hz} $\to$ \omterm{prop:b1:operational-amplifier:schmitt}{Schmitt trigger} at
$\qty{2.0}{V} \pm \qty{0.1}{V}$ $\to$ LED. The real limits: the \omterm{rem:b1:operational-amplifier:real}{input
offset voltage} and the matching of the \omterm{prop:b1:operational-amplifier:sumdiff}{difference amplifier}'s
\omterm{def:b1:dc-circuits:resistor}{resistors}, both comparable to the millivolt \omterm{def:b1:wave-propagation:signal}{signal}.
\end{solution}
